% 第三批：《大学物理简明教程》第5版，3.2.1--3.2.12。
% 源书印刷页69--70，用户扫描PDF物理页同号；质量状态待独立复核。
\subsection{填空题：定轴转动、角动量和能量}

\sourceproblem{3.2.1}{69}{69}
\problemtext 三个各为$m$的质点位于边长$a$的等边三角形顶点。求过重心垂直三角形的轴、过重心且平行任一边的轴、过重心及一个顶点的面内轴的转动惯量。
\workings
顶点到重心的距离均为$a/\sqrt3$，故法向轴
\[
\boxed{J_0=3m(a/\sqrt3)^2=ma^2}.
\]
由平面薄层的垂直轴定理及等边三角形的$120^\circ$旋转对称性，任何一条面内重心轴均有
\[
\boxed{J_1=J_2=\tfrac12J_0=\tfrac12ma^2}.
\]
\crosscheck $J_0=J_x+J_y$；每一质量点的平方距离非负，且两条正交面内重心轴由于三重对称性具有相同转动惯量。

\sourceproblem{3.2.2}{69}{69}
\problemtext 两个厚度、总质量相等而均匀密度不同的圆盘，$\rho_A>\rho_B$；比较绕各自盘心垂直盘面的轴的转动惯量$J_A,J_B$。
\workings
$M=\rho\pi R^2h$；在$M$与厚度$h$固定下，
$R^2=M/(\rho\pi h)$。均匀圆盘$J=\frac12MR^2$，因此
\[
\frac{J_A}{J_B}=\frac{\rho_B}{\rho_A}<1,
\quad \boxed{J_A<J_B}.
\]
\crosscheck 密度较大、质量相同时半径更小，质量更靠近转轴。

\sourceproblem{3.2.3}{69}{69}
\problemtext 刚体绕定轴转动，$J=3.0\,\mathrm{kg\,m^2}$、$\omega_0=6.0\,\mathrm{rad/s}$，受恒阻力矩$M=-12\,\mathrm{N\,m}$，求减速到$2.0\,\mathrm{rad/s}$时转过的角度。
\workings
角加速度$\alpha=M/J=-4\,\mathrm{rad/s^2}$，
耗时$t=(2-6)/(-4)=1\,\mathrm s$，故
\[
\boxed{\theta=\omega_0t+\tfrac12\alpha t^2
=6-2=4\,\mathrm{rad}}.
\]
\crosscheck 也可由$\omega^2-\omega_0^2=2\alpha\theta$得
$(4-36)/(-8)=4$。

\sourceproblem{3.2.4}{69}{69}
\problemtext 两名各为$70\,\mathrm{kg}$的滑冰者，沿相距$10\,\mathrm m$的平行滑行路线以$6.5\,\mathrm{m/s}$相向经过；在交错瞬间抓住一条$10\,\mathrm m$绳的两端，随即绕绳中点转动。求各自对中心的角动量；若把绳缩短到$5\,\mathrm m$，求各自速率。
\workings
对称情形，初始每人距质心$r_0=5$ m，速度垂直半径，
\[
\boxed{L_0=mr_0v_0=70(5)(6.5)=2275\,\mathrm{kg\,m^2/s}}.
\]
忽略冰面摩擦，绳张力沿中心方向，不提供对中心的力矩。
最终$r_1=2.5$ m，故
\[
\boxed{v_1=L_0/(mr_1)=13\,\mathrm{m/s}}.
\]
\crosscheck 距中心减半，速率加倍；缩绳者需做功，故动能不守恒。

\sourceproblem{3.2.5}{70}{70}
\problemtext 子弹沿水平方向射入下端系有质点球、可绕顶端水平轴转动的轻杆/球系统并滞留其中；判断射入的极短时间内与其后系统升高时各有什么守恒量。
\workings
极短碰撞中，轴对转轴的力矩为零，重力矩冲量可忽略，
所以\textbf{关于固定转轴的角动量守恒}：
\[
L_{\rm before}=L_{\rm after}.
\]
碰撞为完全非弹性，\emph{机械能不守恒}。
碰后上摆若忽略轴摩擦和空气阻力，并把地球引力势能计入，
\[
\boxed{K_{\rm after}+U_{\rm after}=K_{\rm top}+U_{\rm top}},
\]
即\textbf{机械能守恒}。
\crosscheck 碰撞阶段角动量守恒不推出转动动能守恒；提升阶段则需要将地球计入保守系统。

\sourceproblem{3.2.6}{70}{70}
\problemtext 质点系合外力恒为零时，判断合外力矩是否必为零，以及动量、角动量、机械能中哪一量必守恒。
\workings
合外力为零满足$\dot{\bm P}=0$，故
\[
\boxed{\bm P=\mathrm{const}}.
\]
但合外力矩可不为零：一对作用点不同、方向相反的外力构成力偶，
即$\sum\bm F=0$而$\sum\bm M\ne0$。故角动量不一定守恒；内摩擦可改变机械能。
\crosscheck 这里无额外“全部内力保守”条件，所以不能推出机械能守恒。

\sourceproblem{3.2.7}{70}{70}
\problemtext 花样滑冰运动员初始转动惯量$J_0$、角速度$\omega_0$；收臂后转动惯量降为原来的$1/3$。求最终角速度。
\workings
假设外力对竖直轴的力矩忽略，
$J_0\omega_0=(J_0/3)\omega_1$，故
\[
\boxed{\omega_1=3\omega_0}.
\]
\crosscheck 转动动能变为原来的三倍；收臂需人体肌肉做功，不能把动能也当作守恒量。

\sourceproblem{3.2.8}{70}{70}
\problemtext 飞轮以角速度$300\,\mathrm{rad/min}$旋转，转动惯量$J=5\,\mathrm{kg\,m^2}$，恒定阻力矩使其在$2$ s内停止。求阻力矩大小。
\workings
$\omega_0=300/60=5\,\mathrm{rad/s}$，
$\alpha=(0-5)/2=-2.5\,\mathrm{rad/s^2}$，故
\[
\boxed{|M|=J|\alpha|=12.5\,\mathrm{N\,m}}.
\]
\crosscheck 原题单位为$\mathrm{rad/min}$而不是每分钟转数；若误当$300\,\mathrm{r/min}$将引入额外$2\pi$因子。

\sourceproblem{3.2.9}{70}{70}
\problemtext 质量$m$的质点位置$\bm r=a\cos(\omega t)\bm i+
b\sin(\omega t)\bm j$，其中$a,b,\omega$均为正常量。求对坐标原点的合外力矩。
\workings
微分两次得$\bm a=-\omega^2\bm r$，故
$\bm F=-m\omega^2\bm r$为指向原点的中心力。
\[
\boxed{\bm M_O=\bm r\times\bm F=\bm0}.
\]
\crosscheck 即使轨迹为椭圆、速率不恒定，中心力对其中心的力矩仍为零。

\sourceproblem{3.2.10}{70}{70}
\problemtext 质量$m$、长度$L$的均匀细杆，可绕其\emph{中点}$O$处垂直于竖直面的水平定轴转动，两端分别固定$m$和$2m$的质点球。系统从与水平成$60^\circ$的静止位置释放。问杆转到水平时重力对系统关于$O$的合力矩大小。
\workings
杆的质心在$O$，故杆重力矩为零。两个球关于$O$的力臂均为$L/2$，但方向相反，合力矩
\[
\boxed{|M_O|=2mg\frac L2-mg\frac L2=\frac12mgL}.
\]
\crosscheck 此结论依“过中点$O$”与原图两端各一个质点球；若轴位置变化，力矩亦变化。

\sourceproblem{3.2.11}{70}{70}
\problemtext 刚体由$\omega_0=8\,\mathrm{rad/s}$在大小$8\,\mathrm{N\,m}$的恒定阻力矩作用下，$2$ s后降至$4\,\mathrm{rad/s}$。求角加速度及转动惯量。
\workings
\[
\boxed{\alpha=\frac{4-8}{2}=-2\,\mathrm{rad/s^2}},
\qquad
\boxed{J=\frac{|M|}{|\alpha|}=4\,\mathrm{kg\,m^2}}.
\]
\crosscheck 阻力矩为负，角加速度也为负，$J$必须为正。

\sourceproblem{3.2.12}{70}{70}
\problemtext 人站在转台中心，两臂侧平举时系统以$2\pi\,\mathrm{rad/s}$旋转，转动惯量为$6\,\mathrm{kg\,m^2}$；收臂后转动惯量为$2\,\mathrm{kg\,m^2}$。求末初转动动能之比。
\workings
关于竖直轴角动量守恒，$\omega_1=(6/2)\omega_0=3\omega_0$，
\[
\boxed{\frac{K_1}{K_0}
=\frac{\tfrac12J_1\omega_1^2}{\tfrac12J_0\omega_0^2}
=\frac{J_0}{J_1}=3}.
\]
\crosscheck 转动能增加三倍，由人体内部做功提供，而非外力矩。
